% Part: sets-functions-relations
% Chapter: sets
% Section: subsets

\documentclass[../../../include/open-logic-section]{subfiles}

\begin{document}

\olfileid[pt]{sfr}{set}{sub}
\olsection{Subconjuntos e conjuntos potência}

\begin{explain}
Com frequência quereremos comparar conjuntos. E um tipo óbvio de comparação
que se pode fazer é o seguinte: \emph{tudo o que está num conjunto está também
no outro}. Esta situação é suficientemente importante para introduzirmos nova
notação.
\end{explain}

\begin{defn}[Subconjunto]
Se todo !!{element} de um conjunto $A$ é também !!a{element} de~$B$, então
dizemos que $A$ é um \emph{subconjunto} de~$B$, e escrevemos $A \subseteq B$. Se
$A$ não é subconjunto de~$B$, escrevemos $A \not\subseteq B$.
Se $A \subseteq B$ mas $A \neq B$, escrevemos $A \subsetneq B$ e dizemos
que $A$ é um \emph{subconjunto próprio} de $B$.
\end{defn}

\begin{ex}
Todo conjunto é subconjunto de si próprio, e $\emptyset$ é subconjunto de todo
conjunto. O conjunto dos números naturais pares é subconjunto do conjunto dos
números naturais. Também $\{ a, b \} \subseteq \{ a, b, c \}$. Mas $\{ a, b, e
\}$ não é subconjunto de $\{ a, b, c \}$.
\end{ex}

\begin{ex}
O número $2$ é um !!{element} do conjunto dos inteiros, ao passo que o
conjunto dos números pares é subconjunto do conjunto dos inteiros. Contudo, um
conjunto pode tanto ser !!a{element} quanto subconjunto de algum outro
conjunto, por exemplo, $\{0\} \in \{0, \{0\}\}$ e também $\{0\} \subseteq \{0,
\{0\}\}$.
\end{ex}

A extensionalidade dá um critério de identidade para conjuntos: $A = B$ se e só se
todo !!{element} de~$A$ é também !!a{element} de~$B$ e vice-versa.
A definição de ``subconjunto'' define $A \subseteq B$ precisamente como a
primeira metade deste critério: todo !!{element} de~$A$ é também
!!a{element} de~$B$. Claro que a definição também se aplica se trocarmos
$A$ e $B$: isto é, $B \subseteq A$ se e só se todo !!{element}
de~$B$ é também !!a{element} de~$A$. E isso, por sua vez, é exatamente a
parte ``vice-versa'' da extensionalidade. Em outras palavras, a extensionalidade
implica que conjuntos são iguais se e só se são subconjuntos um do outro.

\begin{prop}
$A = B$ se e só se $A \subseteq B$ e $B \subseteq A$.
\end{prop}

Agora é também boa ocasião para introduzir mais alguma
notação útil. Ao definir quando $A$ é subconjunto de~$B$ dissemos que
``todo !!{element} de~$A$ é \dots,'' e preenchemos os ``$\dots$'' com
``!!a{element} de $B$''. Mas esta é uma \emph{forma} tão comum de
expressão que convém introduzir notação formal para ela.

\begin{defn}\ollabel{forallxina}
$(\forall x \in A)\phi$ abrevia $\forall x(x \in A \lif
\phi)$. Do mesmo modo, $(\exists x \in A)\phi$ abrevia $\exists x(x
\in A \land \phi)$. 
\end{defn}

Com esta notação, podemos dizer que $A \subseteq B$ se e só se $(\forall
x \in A)x \in B$. 

Passamos agora a considerar um certo tipo de conjunto: o conjunto de todos os
subconjuntos de um dado conjunto. 

\begin{defn}[Conjunto potência]
O conjunto constituído por todos os subconjuntos de um conjunto~$A$ chama-se o
\emph{conjunto potência} de~$A$, e denota-se $\Pow{A}$.
  \[
    \Pow{A} = \Setabs{B}{B \subseteq A} 
  \]
\end{defn}

\begin{ex}
Quais são todos os subconjuntos possíveis de $\{ a, b, c \}$? São:
$\emptyset$, $\{a \}$, $\{b\}$, $\{c\}$, $\{a, b\}$, $\{a, c\}$, $\{b,
c\}$, $\{a, b, c\}$. O conjunto de todos estes subconjuntos é
$\Pow{\{a,b,c\}}$:
\[
\Pow{\{ a, b, c \}} = \{\emptyset, \{a \}, \{b\}, \{c\}, \{a, b\},
\{b, c\}, \{a, c\}, \{a, b, c\}\}
\]
\end{ex}

\begin{prob}
Liste todos os subconjuntos de $\{a, b, c, d\}$.
\end{prob}

\begin{prob}
Mostre que se $A$ tem $n$ !!{element}s, então $\Pow{A}$ tem $2^n$
!!{element}s.
\end{prob}

\end{document}
